\chapter{Tree Techniques}\label{ch:trees}

A tree is a connected graph with $n$ vertices and $n - 1$ edges; equivalently, a graph in which every two vertices are joined
by exactly one simple path. That uniqueness makes trees the friendliest of graphs. Recursion over subtrees is natural, every
path decomposes at a lowest common ancestor, and dynamic programming works cleanly because subtrees do not overlap. This
chapter collects the techniques that turn those facts into algorithms: traversal orders, diameters and farthest vertices,
flattening a tree with an Euler tour, binary lifting and the many ways to compute lowest common ancestors, DP on trees for one
root and for all roots at once, merging DP tables and sets efficiently, and deciding whether two trees have the same shape.

\begin{chapterintro}{\silver\ \gold\ \plat}{Chapter~\ref{ch:graphs} (BFS and DFS) and Chapter~\ref{ch:dsu} (DSU).}
  \item root a tree and compute parents, depths, subtree sizes and orders without deep recursion;
  \item find diameters and the farthest vertex from every vertex;
  \item answer subtree and path queries by flattening the tree with an Euler tour;
  \item compute lowest common ancestors by binary lifting, and offline with Tarjan's algorithm;
  \item design tree DPs, reroot them to get every root's answer, and merge child tables in $\Oh{n^2}$ total;
  \item merge sets small to large, and test trees for isomorphism with canonical labels.
\end{chapterintro}

% ======================================================================================================================
\section[Rooted trees]{Rooted trees and traversal orders \hfill \silver}\label{sec:rooted}

A few counting facts are used constantly. A tree on $n$ vertices has exactly $n - 1$ edges; removing any edge splits it into two
trees, and adding any edge creates exactly one cycle. The degrees sum to $2n - 2$, so the average degree is below 2 and every tree
with at least two vertices has at least two leaves. Most tree problems give the tree as $n - 1$ undirected edges; some give a
parent array instead.

Choosing a root gives every vertex a \defn{parent} (except the root), \defn{children}, a \defn{depth} (the number of edges to the
root) and a \defn{subtree} (itself and all its descendants), as in Figure~\ref{fig:rooted}. Every vertex other than the root has
exactly one parent, so the parent array alone describes the rooted tree, and one traversal computes all of these quantities.

\begin{figure}[htbp]
\centering
\begin{minipage}[c]{6cm}\centering
\begin{forest}
  cptree, for tree={s sep=7mm}
  [0, hot
    [1, cool [3, cool] [4, cool [7, cool]]]
    [2 [5] [6]]]
\end{forest}
\end{minipage}
\hspace{8mm}
\begin{minipage}[c]{5.8cm}
\raggedright\small root 0 has depth 0; $\mathit{parent}(4) = 1$;\\
$\mathit{children}(1) = \{3, 4\}$;\\
$\mathit{subtree}(1) = \{1, 3, 4, 7\}$, size 4 (blue);\\
$\mathit{depth}(7) = 3$; the leaves are 3, 7, 5, 6.\\[4pt]
BFS order: $0, 1, 2, 3, 4, 5, 6, 7$\\
DFS order: $0, 1, 3, 4, 7, 2, 5, 6$
\end{minipage}
\caption{A rooted tree and its basic quantities. In the DFS order, the subtree of 1 occupies the consecutive positions 1--4.}
\label{fig:rooted}
\end{figure}

\subsection{The BFS-order idiom}

Two orders of the vertices are especially useful. In a \emph{BFS order} every parent comes before its children, so reading the
order backwards processes all children before their parent: exactly what a bottom-up computation needs. Building the order needs no
queue beyond the order itself, and no recursion:

\begin{code}
vi order{1}, par(n + 1, 0), depth(n + 1, 0);
rep(i, n) for (int v : adj[order[i]])  // order grows while it is scanned
    if (v != par[order[i]]) {
        par[v] = order[i], depth[v] = depth[order[i]] + 1;
        order.push_back(v);
    }
vi sub(n + 1, 1);  // subtree sizes
per(i, 1, n) sub[par[order[i]]] += sub[order[i]];  // children first
\end{code}

Every vertex is appended once, when its parent is scanned, so this is $\Oh{n}$. Skipping the parent is enough to avoid visiting a
vertex twice, because in a tree no other edge leads back to an already visited vertex. The last loop is the model for bottom-up tree computations in this
chapter: walk the order backwards and let every vertex hand its value to its parent. The complete program for subtree sizes is the
worked problem Subordinates (Chapter~\ref{ch:worked}).

\begin{pitfall}[Recursion depth and parent arrays]
A tree with $2 \cdot 10^5$ vertices may be a single path, and a recursive DFS on it is $2 \cdot 10^5$ calls deep. Most judges
configure a large stack, but a default 1\,MB stack, as on Windows, overflows somewhere between $10^4$ and $10^5$ frames,
depending on the frame size. The BFS order and the explicit-stack DFS of Section~\ref{sec:eulertour} avoid the problem entirely. When the input is a parent array,
check whether the statement guarantees $p_i < i$: if it does, the order $1, 2, \dots, n$ already has parents before children;
if it does not (as in Company Queries~II below), build the children lists and a BFS order first.
\end{pitfall}

\subsection{DFS orders}

A DFS visits a vertex, then each child's whole subtree in turn, so it produces three useful sequences (Figure~\ref{fig:dfsorders}).
The \defn{preorder} lists the vertices when they are entered; every subtree is a contiguous block that starts with its root, which
is the basis of the Euler tour technique (Section~\ref{sec:eulertour}). The \defn{postorder} lists them when they are left; subtrees
are again contiguous, now ending with their root, and children always precede parents. The full \defn{Euler tour} writes a vertex
down every time the walk passes through it, $2n - 1$ entries in all; between the first visits of $u$ and $v$ the walk climbs no
higher than their lowest common ancestor and does pass through it, so the LCA is the shallowest vertex of that stretch, a range
minimum query (Section~\ref{sec:lca}).

\begin{figure}[htbp]
\centering
\begin{tikzpicture}
  \node[vtx] (w0) at (1.2, 0) {0};
  \node[vtx, hot] (w1) at (0.45, -1.05) {1};
  \node[vtx] (w2) at (1.95, -1.05) {2};
  \node[vtx, hot] (w3) at (0, -2.1) {3};
  \node[vtx, hot] (w4) at (0.9, -2.1) {4};
  \node[vtx] (w5) at (1.95, -2.1) {5};
  \draw[ed] (w0) -- (w1) (w0) -- (w2) (w1) -- (w3) (w1) -- (w4) (w2) -- (w5);
  \node[blbl, anchor=east] at (5.3, -0.2) {preorder};
  \arrs{A}{5.45}{-0.2}{0/n, 1/hot, 3/hot, 4/hot, 2/n, 5/n}
  \node[blbl, anchor=east] at (5.3, -1.1) {postorder};
  \arrs{B}{5.45}{-1.1}{3/hot, 4/hot, 1/hot, 5/n, 2/n, 0/n}
  \node[blbl, anchor=east] at (5.3, -2.0) {Euler tour};
  \arrs{C}{5.45}{-2.0}{0/n, 1/n, 3/cool, 1/good, 4/cool, 1/n, 0/n, 2/n, 5/n, 2/n, 0/n}
  \spanbelow{C-3}{C-5}{from the first 3 to the first 4: the shallowest is $1 = \mathrm{lca}(3, 4)$}
\end{tikzpicture}
\caption{The three DFS orders of one tree. In preorder and postorder the subtree of 1 (red) is a contiguous block. In the Euler
tour, the shallowest vertex between the first visits of 3 and 4 is their LCA (green).}
\label{fig:dfsorders}
\end{figure}

\usetemplate{Tree DFS (parent, depth, subtree size, Euler tour)}{Ctrl+Alt+T, E}{cptree}

% ======================================================================================================================
\section[Diameter and farthest vertices]{Diameter and farthest vertices \hfill \gold}\label{sec:diameter}

The \defn{diameter} of a tree is its longest path. Two BFS passes find it: from any vertex, the farthest vertex $a$ is an endpoint
of some diameter, and the farthest vertex from $a$ is the other endpoint $b$ (the worked problem Tree Diameter).

\emph{Why the first BFS lands on a diameter endpoint.} Let $x$--$y$ be a diameter, let $a$ be the farthest vertex from the start
$s$, and let $c$ be the vertex of the path $x$--$y$ closest to $s$; every path from $s$ into $x$--$y$ enters it at $c$. Let $w$ be
the last vertex that the path $s \to a$ shares with the path $s \to c$.
\begin{itemize}
  \item If $w \ne c$, the branch to $a$ leaves before reaching $c$. Since $a$ is at least as far from $s$ as $x$ is,
        $d(w, a) \ge d(w, c) + d(c, x)$, so the path $a \to w \to c \to y$ has length at least $d(c, x) + d(c, y) = d(x, y)$.
  \item If $w = c$, then $d(c, a) \ge d(c, x)$ and $d(c, a) \ge d(c, y)$ for the same reason. The branch to $a$ leaves $c$ towards
        at most one of the two ends; walking from $a$ to $c$ and on to the other end gives a path of length at least $d(x, y)$.
\end{itemize}
Either way a longest path ends at $a$ (Figure~\ref{fig:diameter}). The argument uses only that the graph is a tree and that the
edge weights are non-negative; it fails on graphs with cycles and with negative weights.

\begin{figure}[htbp]
\centering
\tikzset{dlab/.style={font=\sffamily\scriptsize\bfseries, text=fgblue, inner sep=1pt}}
% #1: name/style/text/distance/label side for the eight vertices; #2: drawn under the edges
\newcommand{\diampanel}[2]{\begin{tikzpicture}
  \foreach \nm/\px/\py in {s/0/0, m/0.9/0, m2/1.7/0.45, a/2.5/0.9, x/1.7/-0.6, p/-0.85/0.5, b/-1.65/1.0, q/-0.85/-0.6}
    \coordinate (D\nm) at (\px, \py);
  #2
  \draw[ed] (Ds) -- (Dm) -- (Dm2) -- (Da) (Dm) -- (Dx) (Ds) -- (Dp) -- (Db) (Ds) -- (Dq);
  \foreach \nm/\sy/\txt/\dv/\an in {#1} \node[vtxs, \sy, label={[dlab]\an:\dv}] at (D\nm) {\txt};
\end{tikzpicture}}
\newcommand{\diampath}{\draw[line width=4.5pt, draw=fgverm, opacity=0.28] (Da) -- (Dm2) -- (Dm) -- (Ds) -- (Dp) -- (Db);}
\begin{filmstrip}[3]
  \begin{stepbox}{1. BFS from $s$: farthest is $a$, at 3}
    \diampanel{s/cool/s/0/above, m/n//1/above, m2/n//2/above, a/hot/a/3/above, x/n//2/below, p/n//1/above,
      b/n/b/2/above, q/n//1/below}{}
  \end{stepbox}
  \begin{stepbox}{2. BFS from $a$: farthest is $b$, at 5}
    \diampanel{s/n/s/3/above, m/n//2/above, m2/n//1/above, a/cool/a/0/above, x/n//3/below, p/n//4/above,
      b/hot/b/5/above, q/n//4/below}{\diampath}
  \end{stepbox}
  \begin{stepbox}{3. farthest distance $= \max(d_a, d_b)$}
    \diampanel{s/n/s/3/above, m/n//3/above, m2/n//4/above, a/good/a/5/above, x/n//4/below, p/n//4/above,
      b/good/b/5/above, q/n//4/below}{\diampath}
  \end{stepbox}
\end{filmstrip}
\caption{The two-BFS diameter algorithm; the numbers are BFS distances. The farthest vertex from $s$ is $a$, the farthest from $a$
is $b$, and $a$--$b$ (red) is a diameter. Panel~3 is the corollary: the farthest vertex from any vertex is $a$ or $b$.}
\label{fig:diameter}
\end{figure}

The same exchange argument proves a stronger fact: for \emph{every} vertex $v$, a farthest vertex from $v$ is one of the two ends
$a, b$ of a fixed diameter. So after computing distances from $a$ and from $b$, the farthest distance from $v$ is $\max(d_a(v),
d_b(v))$.

\begin{problem}{Tree Distances I (CSES)}
For every vertex of a tree, find the largest distance to another vertex.
\end{problem}

\begin{code}
void solve() {
    int n;
    cin >> n;
    vvi adj(n + 1);
    rep(i, n - 1) {
        int a, b;
        cin >> a >> b;
        adj[a].push_back(b);
        adj[b].push_back(a);
    }
    auto bfs = [&](int s) {  // distances from s
        vi d(n + 1, -1);
        vi q{s};
        d[s] = 0;
        rep(i, sz(q)) for (int v : adj[q[i]])
            if (d[v] < 0) d[v] = d[q[i]] + 1, q.push_back(v);
        return d;
    };
    // the farthest vertex from anywhere is an end of a diameter (a or b)
    vi d1 = bfs(1);
    int a = max_element(all(d1)) - d1.begin();
    vi da = bfs(a);
    int b = max_element(all(da)) - da.begin();
    vi db = bfs(b);
    vi ans(n);
    rep(v, 1, n + 1) ans[v - 1] = max(da[v], db[v]);
    print(ans);
}
\end{code}

\usetemplate{Tree diameter}{Ctrl+Alt+T, Shift+D}{cpdiameter}

The \defn{center} of a tree (the middle vertex or middle edge of any diameter) minimises the largest distance to all vertices, and
it is found by peeling leaves layer by layer until one or two vertices remain (used in Section~\ref{sec:treeiso}). The
\defn{centroid}, a vertex whose removal leaves components of size at most $n/2$, is a different notion; it drives the
decompositions of Chapter~\ref{ch:treedecomp}.

% ======================================================================================================================
\section[The Euler tour technique]{The Euler tour technique \hfill \gold}\label{sec:eulertour}

Number the vertices in the order a DFS enters them: $\mathit{tin}(v)$ is the entry time, and $\mathit{tout}(v)$ the time when the
DFS leaves $v$, which is one more than the largest entry time in its subtree. Then the subtree of $v$ is exactly the contiguous
range of positions $[\mathit{tin}(v), \mathit{tout}(v))$ (Figure~\ref{fig:eulertour}). This \defn{Euler tour technique} flattens the
tree into an array, and every subtree operation becomes a range operation that a Fenwick or segment tree supports. As a bonus, $u$
is an ancestor of $v$ exactly when $\mathit{tin}(u) \le \mathit{tin}(v) < \mathit{tout}(u)$, an $\Oh{1}$ ancestor test.

\begin{figure}[htbp]
\centering
\begin{tikzpicture}
  \tikzset{tl/.style={font=\sffamily\scriptsize, text=fgblue, inner sep=1.5pt}}
  \node[vtx] (v0) at (1.4, 0) {0};
  \node[vtx, hot] (v1) at (0.5, -1.1) {1};
  \node[vtx] (v2) at (2.4, -1.1) {2};
  \node[vtx, hot] (v3) at (0, -2.2) {3};
  \node[vtx, hot] (v4) at (1.0, -2.2) {4};
  \node[vtx] (v5) at (2.4, -2.2) {5};
  \draw[ed] (v0) -- (v1) (v0) -- (v2) (v1) -- (v3) (v1) -- (v4) (v2) -- (v5);
  \node[tl, anchor=west] at (v0.east) {$[0, 6)$};
  \node[tl, anchor=east] at (v1.west) {$[1, 4)$};
  \node[tl, anchor=west] at (v2.east) {$[4, 6)$};
  \node[tl, anchor=north] at (v3.south) {$[2, 3)$};
  \node[tl, anchor=north] at (v4.south) {$[3, 4)$};
  \node[tl, anchor=north] at (v5.south) {$[5, 6)$};
  \node[blbl, anchor=east] at (5.5, -1.1) {vertex};
  \arrs{E}{5.7}{-1.1}{0/n, 1/hot, 3/hot, 4/hot, 2/n, 5/n}
  \idx{E}{6}
  \node[lbl, anchor=east] at (5.5, -1.63) {position};
  \spanabove{E-2}{E-4}{subtree of 1}
\end{tikzpicture}
\caption{The Euler tour technique. Each vertex shows its range $[\mathit{tin}, \mathit{tout})$; listing the vertices by entry time
makes every subtree a contiguous block, here positions $[1, 4)$ for the subtree of 1.}
\label{fig:eulertour}
\end{figure}

\begin{problem}{Subtree Queries (CSES)}
Every vertex of a rooted tree has a value. Process updates ``set the value of $s$ to $x$'' and queries ``the sum of the values in
the subtree of $s$''.
\solution Flatten with an iterative DFS and keep the values in a Fenwick tree at position $\mathit{tin}(v)$. A subtree sum is the sum
over positions $[\mathit{tin}(s), \mathit{tout}(s))$, and an update adds the difference at one position.
\end{problem}

\begin{code}
void solve() {  // BIT (cpbit) is pasted above solve()
    int n, q;
    cin >> n >> q;
    vi val(n + 1);
    rep(v, 1, n + 1) cin >> val[v];
    vvi adj(n + 1);
    rep(i, n - 1) {
        int a, b;
        cin >> a >> b;
        adj[a].push_back(b);
        adj[b].push_back(a);
    }
    // Euler tour: the subtree of v is the position range [tin[v], tout[v])
    vi tin(n + 1), tout(n + 1), st{1}, par(n + 1, 0), it(n + 1, 0);
    int timer = 0;
    tin[1] = timer++;
    while (!st.empty()) {  // iterative DFS
        int u = st.back();
        if (it[u] < sz(adj[u])) {
            int v = adj[u][it[u]++];
            if (v == par[u]) continue;
            par[v] = u, tin[v] = timer++;
            st.push_back(v);
        } else {
            tout[u] = timer;
            st.pop_back();
        }
    }
    BIT bit(n);
    rep(v, 1, n + 1) bit.add(tin[v], val[v]);
    while (q--) {
        int type, s;
        cin >> type >> s;
        if (type == 1) {
            ll x;
            cin >> x;
            bit.add(tin[s], x - val[s]), val[s] = x;
        } else print(bit.sum(tin[s], tout[s]));  // the positions of s's subtree
    }
}
\end{code}

\subsection{Paths to the root}

The same flattening handles sums along the path from the root to a vertex, with one twist. The value of $v$ contributes to the path
sum of every vertex in $v$'s subtree. So add it at position $\mathit{tin}(v)$ and subtract it at $\mathit{tout}(v)$: the prefix sum
up to $\mathit{tin}(u)$ then contains $v$'s value exactly when $u$ lies in $v$'s subtree, that is, when $v$ is on the path from the
root to $u$ (Figure~\ref{fig:pathsum}). A path between any $u$ and $w$ follows by inclusion--exclusion: $\mathit{path}(u) +
\mathit{path}(w) - 2\,\mathit{path}(\ell) + \mathit{val}(\ell)$ with $\ell = \mathrm{lca}(u, w)$.

\begin{figure}[htbp]
\centering
\begin{tikzpicture}
  \node[blbl, anchor=east] at (-0.2, 0.64) {vertex};
  \foreach \vv/\cl [count=\i] in {0/fgink, 1/fgblue, 3/fgblue, 4/fgblue, 2/fgink, 5/fgink}
    \node[font=\sffamily\small\bfseries, text=\cl] at ({\i*0.7 - 0.35}, 0.64) {\vv};
  \node[blbl, anchor=east] at (-0.2, 0) {added};
  \arrs{P}{0}{0}{0/n, $+x$/good, 0/n, 0/n, $-x$/hot, 0/n, 0/n}
  \idx{P}{7}
  \node[blbl, anchor=east] at (-0.2, -1.45) {prefix sum};
  \arrs{Q}{0}{-1.45}{0/n, $x$/cool, $x$/cool, $x$/cool, 0/n, 0/n, 0/n}
  \node[notel, anchor=west, text width=5.6cm] at (5.3, -0.7) {Vertex 1 has value $x$. It is added at $\mathit{tin}(1) = 1$
    and removed at $\mathit{tout}(1) = 4$, so the prefix sums at positions 1, 2 and 3 (the vertices 1, 3, 4 of its subtree)
    include it, and no other prefix sum does.};
\end{tikzpicture}
\caption{Root-path sums with an Euler tour, for the value $x$ of vertex 1 alone in the tree of Figure~\ref{fig:eulertour}. A value
is added where its subtree's range starts and removed where it ends.}
\label{fig:pathsum}
\end{figure}

\begin{problem}{Path Queries (CSES)}
Every vertex has a value. Process updates of single values and queries for the sum of the values on the path from the root to a
vertex.
\end{problem}

\begin{code}
void solve() {  // BIT (cpbit) is pasted above solve()
    int n, q;
    cin >> n >> q;
    vi val(n + 1);
    rep(v, 1, n + 1) cin >> val[v];
    vvi adj(n + 1);
    rep(i, n - 1) {
        int a, b;
        cin >> a >> b;
        adj[a].push_back(b);
        adj[b].push_back(a);
    }
    vi tin(n + 1), tout(n + 1), st{1}, par(n + 1, 0), it(n + 1, 0);
    int timer = 0;
    tin[1] = timer++;
    while (!st.empty()) {
        int u = st.back();
        if (it[u] < sz(adj[u])) {
            int v = adj[u][it[u]++];
            if (v == par[u]) continue;
            par[v] = u, tin[v] = timer++;
            st.push_back(v);
        } else {
            tout[u] = timer;
            st.pop_back();
        }
    }
    // a value at v counts for every vertex in v's subtree: add it at tin[v]
    // and remove it at tout[v]; a prefix sum then gives the root path sum
    BIT bit(n + 1);
    auto change = [&](int v, ll d) {
        bit.add(tin[v], d), bit.add(tout[v], -d);
    };
    rep(v, 1, n + 1) change(v, val[v]);
    while (q--) {
        int type, s;
        cin >> type >> s;
        if (type == 1) {
            ll x;
            cin >> x;
            change(s, x - val[s]), val[s] = x;
        } else print(bit.sum(0, tin[s] + 1));  // positions 0..tin[s]
    }
}
\end{code}

\begin{recognize}
``Update a vertex, query a subtree'', ``update a subtree, query a vertex'' (a range add and a point query), ``path from the
root'' sums, and ancestor tests. When both endpoints of a path vary and values change, heavy-light decomposition takes over
(Chapter~\ref{ch:treedecomp}).
\end{recognize}

% ======================================================================================================================
\section[Binary lifting and the LCA]{Binary lifting and the lowest common ancestor \hfill \gold}\label{sec:lca}

The \defn{lowest common ancestor} $\mathrm{lca}(u, v)$ is the deepest vertex that is an ancestor of both. The path between $u$ and
$v$ climbs from $u$ to the LCA and descends to $v$, so $\mathrm{dist}(u, v) = \mathit{depth}(u) + \mathit{depth}(v) -
2\,\mathit{depth}(\mathrm{lca})$, and many path questions reduce to it.

\defn{Binary lifting} stores $\mathit{up}[j][v]$, the ancestor $2^j$ levels above $v$, computed level by level with
$\mathit{up}[j][v] = \mathit{up}[j-1][\,\mathit{up}[j-1][v]\,]$ in $\Oh{n \log n}$. A jump of $k$ levels is then a sum of powers of
two: one lookup per set bit of $k$ (Figure~\ref{fig:lifting}).

\begin{figure}[htbp]
\centering
\begin{tikzpicture}
  \tikzset{jl/.style={font=\sffamily\scriptsize\bfseries, inner sep=1.5pt, text=fgverm}}
  \foreach \dd/\sy in {0/n, 1/n, 2/n, 3/good, 4/n, 5/n, 6/n, 7/n, 8/cool}
    \node[vtxs, \sy] (p\dd) at ({0.35 + 0.7*\dd}, 0.95) {\dd};
  \foreach \dd in {1,...,8} { \pgfmathtruncatemacro{\uu}{\dd - 1} \draw[ed] (p\uu) -- (p\dd); }
  \draw[->, line width=1.1pt, draw=fgverm] (p8) to[bend right=45] node[jl, above] {$2^2$} (p4);
  \draw[->, line width=1.1pt, draw=fgverm] (p4) to[bend right=80] node[jl, above] {$2^0$} (p3);
  \def\dotcell{$\cdot$}
  \arrs{U0}{0}{0}{\dotcell/n, 0/n, 1/n, 2/n, 3/hot, 4/n, 5/n, 6/n, 7/cool}
  \arrs{U1}{0}{-0.7}{\dotcell/n, \dotcell/n, 0/n, 1/n, 2/n, 3/n, 4/n, 5/n, 6/cool}
  \arrs{U2}{0}{-1.4}{\dotcell/n, \dotcell/n, \dotcell/n, \dotcell/n, 0/n, 1/n, 2/n, 3/n, 4/hot}
  \arrs{U3}{0}{-2.1}{\dotcell/n, \dotcell/n, \dotcell/n, \dotcell/n, \dotcell/n, \dotcell/n, \dotcell/n, \dotcell/n, 0/cool}
  \foreach \jj in {0, 1, 2, 3}
    \node[blbl, anchor=east] at ({-0.15}, {-0.7*\jj}) {$\mathit{up}[\jj][v]$};
  \node[notel, anchor=west, text width=5.4cm] at (6.9, -0.55) {Row $j$ holds the ancestor $2^j$ levels up (a dot: none), and
    is built from row $j - 1$, e.g.\ $\mathit{up}[2][8] = \mathit{up}[1][\mathit{up}[1][8]] = \mathit{up}[1][6] = 4$. Column 8
    (blue) is everything vertex 8 stores. Climbing $k = 5 = 4 + 1$ levels takes one lookup per set bit (red):
    $\mathit{up}[2][8] = 4$, then $\mathit{up}[0][4] = 3$.};
\end{tikzpicture}
\caption{Binary lifting on a path of 9 vertices rooted at 0. With $\lceil \log_2 n \rceil$ rows in the table, any ancestor is at
most that many lookups away.}
\label{fig:lifting}
\end{figure}

\begin{problem}{Company Queries I (CSES)}
Given the boss of every employee, answer queries: who is the $k$-th boss above employee $x$ (or $-1$ if there is none)?
\end{problem}

\begin{code}
void solve() {
    int n, q;
    cin >> n >> q;
    const int LOG = 18;  // 2^18 > 2 * 10^5
    // up[j][v] = the boss 2^j levels above v, 0 if there is none
    vector<vi> up(LOG, vi(n + 1, 0));
    rep(v, 2, n + 1) cin >> up[0][v];
    rep(j, 1, LOG) rep(v, 1, n + 1) up[j][v] = up[j - 1][up[j - 1][v]];
    while (q--) {
        int x, k;
        cin >> x >> k;
        // k is a sum of powers of two: one jump per set bit
        for (int j = 0; j < LOG && x; j++)
            if (k >> j & 1) x = up[j][x];
        print(x ? x : -1);
    }
}
\end{code}

The table has $\lceil \log_2 n \rceil$ rows of $n + 1$ entries, $18 \times 2 \cdot 10^5$ numbers here; with \cc{#define int long long}
that is about 29\,MB, fine for a 256\,MB limit but worth a thought at $n = 10^6$, where an \cc{int32_t} table halves it. Keeping
the level $j$ as the \emph{outer} index makes the construction loop scan memory sequentially. Nothing in the construction needs a
tree: the same table answers ``where am I after $k$ steps'' in any \emph{functional graph}, where every vertex has one successor
(Planets Queries~I, Chapter~\ref{ch:worked}).

\usetemplate{K-th successor (binary lifting)}{Ctrl+Alt+T, .}{cpjump}

For the LCA, first lift the deeper vertex to the depth of the other. If they now coincide, that vertex is the LCA. Otherwise, for
$j$ from the largest down, jump both vertices by $2^j$ whenever that keeps them \emph{different}; this is binary search by jumps
(Section~\ref{sec:predicate}) on the predicate ``the ancestors at this height differ''. At the end they are the two children of the
LCA just below it, and the LCA is their common parent (Figure~\ref{fig:lcasteps}).

\begin{figure}[htbp]
\centering
\begin{minipage}[c]{5cm}\centering
\begin{forest}
  cptree, for tree={s sep=6mm, l sep=4mm}
  [1, good [2, warn [4, warn [7 [10, cool]]] [5]] [3, warn [6, warn [8 [9, cool]]]]]
\end{forest}
\end{minipage}
\hspace{8mm}
\begin{minipage}[c]{6.6cm}
\raggedright\small $\mathrm{lca}(10, 9)$, both at depth 4, so no levelling is needed (blue).\\[3pt]
$j = 2$: jumping 4 levels reaches 1 from both: equal, skip.\\
$j = 1$: jumping 2 levels gives 4 and 6: different, take it.\\
$j = 0$: jumping 1 level gives 2 and 3: different, take it.\\[3pt]
The two vertices (orange) are now children of the LCA, so $\mathrm{lca}(10, 9) = \mathit{parent}(2) = 1$ (green).
\end{minipage}
\caption{Computing an LCA by binary lifting: bring both vertices to the same depth, then take the largest jumps that keep them
apart.}
\label{fig:lcasteps}
\end{figure}

\begin{pseudo}{LCA by binary lifting}
\begin{algorithmic}[1]
\If{$\mathit{depth}(a) < \mathit{depth}(b)$} \State swap $a$ and $b$ \EndIf
\For{$j \gets 0$ \textbf{to} $L - 1$} \Comment{climb $\mathit{depth}(a) - \mathit{depth}(b)$ levels}
  \If{bit $j$ of $\mathit{depth}(a) - \mathit{depth}(b)$ is set} \State $a \gets \mathit{up}[j][a]$ \EndIf
\EndFor
\If{$a = b$} \State \Return $a$ \EndIf
\For{$j \gets L - 1$ \textbf{down to} $0$}
  \If{$\mathit{up}[j][a] \ne \mathit{up}[j][b]$} \State $a \gets \mathit{up}[j][a]$; $b \gets \mathit{up}[j][b]$ \EndIf
\EndFor
\State \Return $\mathit{up}[0][a]$
\end{algorithmic}
\end{pseudo}

\emph{Why it is correct.} After levelling, let $h$ be the number of levels between $a$ (or $b$) and their LCA. The ancestors of $a$
and $b$ that are $t$ levels up are different for $t < h$ and equal for $t \ge h$, a monotone predicate in $t$. The downward loop
builds the largest $t$ with different ancestors, $t = h - 1$, bit by bit from the top, exactly like binary search by jumps; one more
step up reaches the LCA. For the root, $\mathit{up}[j]$ points at the root itself, so an overshooting jump lands on equal vertices
and is simply not taken. Each query costs $2\lceil \log_2 n\rceil$ lookups.

\begin{problem}{Company Queries II (CSES)}
Answer queries: the lowest common boss of employees $a$ and $b$.
\end{problem}

\begin{code}
void solve() {
    int n, q;
    cin >> n >> q;
    const int LOG = 18;
    vector<vi> up(LOG, vi(n + 1, 1));  // above the root: the root itself
    vvi kids(n + 1);
    rep(v, 2, n + 1) cin >> up[0][v], kids[up[0][v]].push_back(v);
    vi depth(n + 1, 0), order{1};  // depths in BFS order from the root
    rep(i, n) for (int c : kids[order[i]]) {
        depth[c] = depth[order[i]] + 1;
        order.push_back(c);
    }
    rep(j, 1, LOG) rep(v, 1, n + 1) up[j][v] = up[j - 1][up[j - 1][v]];
    auto lca = [&](int a, int b) {
        if (depth[a] < depth[b]) swap(a, b);
        int diff = depth[a] - depth[b];
        rep(j, LOG) if (diff >> j & 1) a = up[j][a];  // same depth now
        if (a == b) return a;
        per(j, LOG) if (up[j][a] != up[j][b]) a = up[j][a], b = up[j][b];
        return up[0][a];  // largest jumps that keep them apart, then one more
    };
    while (q--) {
        int a, b;
        cin >> a >> b;
        print(lca(a, b));
    }
}
\end{code}

\usetemplate{LCA (binary lifting)}{Ctrl+Alt+T, L}{cplca}

\begin{trick}[Other ways to get an LCA]
\begin{itemize}
  \item \textbf{Euler tour + sparse table}: write down a vertex every time the DFS passes through it (length $2n - 1$); the LCA of
        $u$ and $v$ is the shallowest vertex between their first occurrences, a range minimum: $\Oh{n \log n}$ preprocessing and
        $\Oh{1}$ per query (Chapter~\ref{ch:static}).
  \item \textbf{Heavy-light decomposition}: $\Oh{\log n}$ per query with $\Oh{n}$ memory and a tiny constant
        (Chapter~\ref{ch:treedecomp}).
  \item \textbf{Offline}: Tarjan's algorithm, next.
  \item \textbf{Path aggregates}: store the maximum (or sum, or minimum) of the $2^j$ edges above $v$ next to $\mathit{up}[j][v]$,
        and combine them during the jumps: maximum edge on a path in $\Oh{\log n}$.
\end{itemize}
\end{trick}

\begin{trick}[Jump pointers with one pointer per vertex]
When vertices arrive one at a time as new leaves, and memory is tight, a single extra pointer per vertex still gives
$\Oh{\log n}$ ancestor queries. Each vertex stores its parent and one \emph{jump}, set when the vertex is added:
\begin{code}
void add_leaf(int v, int p) {  // the root has par = jump = itself
    par[v] = p, depth[v] = depth[p] + 1;
    int j = jump[p];
    jump[v] = depth[p] - depth[j] == depth[j] - depth[jump[j]] ? jump[j] : p;
}
int ancestor_at_depth(int v, int d) {  // requires d <= depth[v]
    while (depth[v] > d) v = depth[jump[v]] >= d ? jump[v] : par[v];
    return v;
}
\end{code}
The jump lengths follow the skew-binary number system, so any climb takes at most about $2\log_2 n$ steps (41 on a path of $10^6$
vertices, measured). For the LCA, level the two vertices, then move both by their jumps when the jumps differ and by their
parents otherwise; at equal depths the jumps have equal lengths.
\end{trick}

\usetemplate{Jump pointers (one pointer per vertex)}{Ctrl+Alt+U, J}{cpjumpptr}

% ======================================================================================================================
\section[Offline LCA]{Offline LCA: Tarjan's algorithm \hfill \plat}\label{sec:offlinelca}

When all queries are known in advance, one DFS with a DSU answers them in nearly linear time. When the DFS \emph{finishes} a vertex,
it merges that vertex's set into its parent's set, and the parent becomes the \emph{anchor} of the merged set. At the moment the DFS
finishes $u$, the finished vertices are grouped into sets, one for every open ancestor of $u$, each containing that ancestor's
finished part of the tree (Figure~\ref{fig:tarjanlca}). So for a query $(u, w)$ in which $w$ is already finished, the anchor of $w$'s
set is exactly the deepest open ancestor of $w$, and that is $\mathrm{lca}(u, w)$. Every query is answered when its later endpoint
finishes.

\begin{figure}[htbp]
\centering
\begin{tikzpicture}
  \fill[fgsky!22, rounded corners=8pt] (-2.85, -2.45) rectangle (-0.45, -0.55);
  \fill[fggreen!16, rounded corners=8pt] (0.5, -2.45) rectangle (1.5, -1.55);
  \node[vtx, warn] (r) at (0, 0) {1};
  \node[vtx] (a) at (-1.5, -1.0) {2}; \node[vtx, warn] (b) at (1.0, -1.0) {3};
  \node[vtx] (c) at (-2.3, -2.0) {4}; \node[vtx] (d) at (-0.9, -2.0) {5}; \node[vtx, hot] (e) at (2.3, -2.0) {6};
  \node[vtx] (f) at (1.0, -2.0) {7};
  \draw[ed] (r) -- (a) (r) -- (b) (a) -- (c) (a) -- (d) (b) -- (f) (b) -- (e);
  \node[font=\sffamily\scriptsize\bfseries, text=fgblue] at (-1.65, -2.7) {set anchored at 1};
  \node[font=\sffamily\scriptsize\bfseries, text=fggreen!70!black] at (1.0, -2.7) {anchored at 3};
  \legendat{(3.4, 0.1)}{hot/finishing now, warn/open: on the DFS stack, n/finished}
  \node[notel, anchor=north west, text width=5.6cm] at (3.3, -1.25) {Queries answered now:
    $\mathrm{lca}(6, 5) = 1$, the anchor of the set of 5, and $\mathrm{lca}(6, 7) = 3$, the anchor of the set of 7.};
\end{tikzpicture}
\caption{Tarjan's offline LCA in the middle of the DFS. Every finished vertex belongs to a set (shaded) that hangs from an open
ancestor, its anchor; for a finished $w$, the anchor of $w$'s set is its LCA with the vertex being finished.}
\label{fig:tarjanlca}
\end{figure}

\begin{problem}{Distance Queries (CSES)}
Answer queries for the distance between two vertices of a tree.
\end{problem}

\begin{code}
void solve() {  // DSU (cpdsu) is pasted above solve()
    int n, q;
    cin >> n >> q;
    vvi adj(n + 1);
    rep(i, n - 1) {
        int a, b;
        cin >> a >> b;
        adj[a].push_back(b);
        adj[b].push_back(a);
    }
    vector<vector<pii>> ask(n + 1);  // (other vertex, query index)
    rep(i, q) {
        int a, b;
        cin >> a >> b;
        ask[a].push_back({b, i});
        ask[b].push_back({a, i});
    }
    // Tarjan's offline LCA: a finished vertex joins its parent's set, whose
    // anchor is the parent. For a query (u, w) answered when u finishes and
    // w is already finished, the anchor of w's set is lca(u, w).
    DSU d(n + 1);
    vi anchor(n + 1), depth(n + 1, 0), par(n + 1, 0), it(n + 1, 0), ans(q);
    iota(all(anchor), 0);
    vector<bool> done(n + 1, false);
    vi st{1};
    while (!st.empty()) {  // iterative DFS
        int u = st.back();
        if (it[u] < sz(adj[u])) {
            int v = adj[u][it[u]++];
            if (v == par[u]) continue;
            par[v] = u, depth[v] = depth[u] + 1;
            st.push_back(v);
            continue;
        }
        done[u] = true;
        for (auto [w, i] : ask[u]) {
            if (!done[w]) continue;
            int l = anchor[d.find(w)];  // lca(u, w)
            ans[i] = depth[u] + depth[w] - 2 * depth[l];
        }
        st.pop_back();
        if (par[u]) d.unite(u, par[u]), anchor[d.find(u)] = par[u];
    }
    for (ll a : ans) print(a);
}
\end{code}

\usetemplate{Offline LCA (Tarjan)}{Ctrl+Alt+U, Shift+L}{cpofflinelca}

\emph{Complexity and when to use it.} Each vertex is pushed and popped once, each query is looked at twice (once from each
endpoint; only the second look answers it), and the DSU operations cost $\Oh{\alpha(n)}$ each: $\Oh{(n + q)\,\alpha(n)}$ in
total, with arrays only. That makes it the fastest LCA method in practice when all queries are known in advance, and it needs
$\Oh{n + q}$ memory instead of $\Oh{n \log n}$. A query $(u, u)$ is answered with $\mathrm{lca} = u$, since $u$'s own set is
anchored at $u$ when it finishes. When queries arrive online, or depend on earlier answers, use binary lifting or the Euler tour
with a sparse table instead.

% ======================================================================================================================
\section[DP on trees]{DP on trees \hfill \gold}\label{sec:treedp}

Most optimisation and counting problems on trees are solved by computing a small table for every subtree from its children's tables,
children first. The state is almost always ``the vertex, plus a little information about the vertex itself'' (chosen or not, matched
or not, its colour), and the transition combines the children independently, because their subtrees do not interact except through
the parent.

\paragraph{A first example: maximum-weight independent set.} Choose vertices of maximum total weight so that no two chosen vertices
are adjacent. For each vertex keep two numbers: $\mathit{in}(v)$, the best total in $v$'s subtree when $v$ is chosen, and
$\mathit{out}(v)$, the best when it is not. A chosen vertex forces its children out; an unchosen vertex leaves each child free:
\[
  \mathit{in}(v) = w(v) + \sum_{c} \mathit{out}(c), \qquad \mathit{out}(v) = \sum_{c} \max\bigl(\mathit{in}(c), \mathit{out}(c)\bigr),
\]
summing over the children $c$ of $v$. The answer is $\max(\mathit{in}(r), \mathit{out}(r))$ at the root $r$. Every vertex and every
edge is handled once, so the DP is $\Oh{n}$. Counting versions have the same shape with products in place of sums: the number of
independent sets satisfies $\mathit{in}(v) = \prod_c \mathit{out}(c)$ and $\mathit{out}(v) = \prod_c (\mathit{in}(c) +
\mathit{out}(c))$, because the children's choices are independent.

\paragraph{The recipe.} Root the tree anywhere. Ask what a parent needs to know about a child's subtree to extend a solution (is the
child chosen, matched, coloured red?) and make that the state. Write the step that folds one child into its parent; it must not
depend on the order of the children. Process the vertices children first, in reverse BFS order, and read the answer at the root.

\begin{problem}{Tree Matching (CSES)}
Find a maximum matching in a tree: as many edges as possible, no two sharing a vertex.
\solution For every vertex $v$ keep two values: $\mathit{free}(v)$, the best matching in $v$'s subtree when $v$ itself stays
unmatched, and $\mathit{best}(v)$, the best when $v$ may be matched. Then $\mathit{free}(v) = \sum_c \mathit{best}(c)$ over the
children, and matching $v$ with a child $c$ changes the total by $\mathit{free}(c) + 1 - \mathit{best}(c)$; the best such gain
gives $\mathit{best}(v)$ (Figure~\ref{fig:treematching}). A greedy algorithm, bottom-up, that matches each vertex with its parent
whenever both are still free, is also optimal, by an exchange argument.
\end{problem}

\begin{figure}[htbp]
\centering
\begin{tikzpicture}
  \tikzset{fb/.style={font=\sffamily\scriptsize, text=fgblue, inner sep=2pt}}
  \node[vtx] (m1) at (1.5, 0) {1};
  \node[vtx] (m2) at (0.3, -1.1) {2}; \node[vtx] (m3) at (2.7, -1.1) {3};
  \node[vtx] (m4) at (-0.3, -2.2) {4}; \node[vtx] (m5) at (0.9, -2.2) {5}; \node[vtx] (m6) at (2.7, -2.2) {6};
  \node[vtx] (m7) at (2.7, -3.3) {7};
  \draw[ed] (m2) -- (m5);
  \draw[hed] (m1) -- (m3) (m2) -- (m4) (m6) -- (m7);
  \draw[ed] (m1) -- (m2) (m3) -- (m6);
  \node[fb, anchor=west] at (m1.east) {$(2, 3)$};
  \node[fb, anchor=east] at (m2.west) {$(0, 1)$};
  \node[fb, anchor=west] at (m3.east) {$(1, 1)$};
  \node[fb, anchor=north] at (m4.south) {$(0, 0)$};
  \node[fb, anchor=north] at (m5.south) {$(0, 0)$};
  \node[fb, anchor=west] at (m6.east) {$(0, 1)$};
  \node[fb, anchor=west] at (m7.east) {$(0, 0)$};
  \node[notel, anchor=west, text width=6.4cm] at (4.3, -1.5) {Each vertex shows $(\mathit{free}, \mathit{best})$. At 3:
    $\mathit{free}(3) = \mathit{best}(6) = 1$, and matching 3 with 6 would gain $\mathit{free}(6) + 1 - \mathit{best}(6) = 0$,
    so $\mathit{best}(3) = 1$. At the root: $\mathit{free}(1) = 1 + 1 = 2$, and matching 1 with 3 gains $1 + 1 - 1 = 1$, so the
    maximum matching has $3$ edges (red).};
\end{tikzpicture}
\caption{The tree-matching DP on a small tree, and a maximum matching read back from it.}
\label{fig:treematching}
\end{figure}

\begin{code}
void solve() {
    int n;
    cin >> n;
    vvi adj(n + 1);
    rep(i, n - 1) {
        int a, b;
        cin >> a >> b;
        adj[a].push_back(b);
        adj[b].push_back(a);
    }
    vi order{1}, par(n + 1, 0);
    rep(i, n) for (int v : adj[order[i]])
        if (v != par[order[i]]) par[v] = order[i], order.push_back(v);
    // free[v]: best matching in v's subtree with v left unmatched;
    // best[v]: best matching in v's subtree with v allowed to be matched
    vi freeV(n + 1, 0), best(n + 1, 0);
    per(i, n) {  // children before parents
        int v = order[i];
        ll gain = 0;  // the most we gain by matching v with one child
        for (int c : adj[v]) {
            if (c == par[v]) continue;
            freeV[v] += best[c];
            chmax(gain, freeV[c] + 1 - best[c]);
        }
        best[v] = freeV[v] + gain;
    }
    print(best[1]);
}
\end{code}

\usetemplate{Tree DP (subtree DP, iterative)}{Ctrl+Alt+T, ;}{cptreedp}

\begin{recognize}
An optimum or a count over choices of vertices or edges of a tree: matchings, independent sets, dominating sets, colourings, ``choose
$k$ vertices'', ``count subtrees with a property''. Write the state for a vertex, decide what a child must tell its parent, and
process vertices in reverse BFS order to avoid deep recursion.
\end{recognize}

% ======================================================================================================================
\section[Rerooting]{DP for every root: rerooting \hfill \gold}\label{sec:rerooting}

Some problems want the answer for \emph{every} choice of root, such as the sum of distances from each vertex to all others. Running
the DP from every root costs $\Oh{n^2}$. \defn{Rerooting} computes all answers in two passes. The first pass roots the tree at one
vertex and computes subtree values. The second pass moves the root across each edge, from a parent $p$ to its child $v$, and updates
the answer in $\Oh{1}$ from $p$'s answer: only the edge $p$--$v$ changes sides (Figure~\ref{fig:reroot}).

\begin{figure}[htbp]
\centering
\begin{tikzpicture}
  \fill[fgsky!18, rounded corners=8pt] (0.6, -0.55) rectangle (3.4, -2.75);
  \fill[fgverm!9, rounded corners=8pt] (-1.9, 0.42) rectangle (0.42, -1.38);
  \node[vtx, hot] (p) at (0, 0) {$p$}; \node[vtx, cool] (v) at (2.0, -1.0) {$v$};
  \node[vtx] (x1) at (-1.4, -1.0) {}; \node[vtx] (x2) at (-0.5, -1.0) {};
  \node[vtx] (y1) at (1.3, -2.1) {}; \node[vtx] (y2) at (2.7, -2.1) {};
  \draw[ed] (p) -- (v) (p) -- (x1) (p) -- (x2) (v) -- (y1) (v) -- (y2);
  \node[font=\sffamily\scriptsize, text=fgblue] at (2.0, -3.05) {$\mathit{sub}(v)$ vertices, each 1 closer};
  \node[font=\sffamily\scriptsize, text=fgverm] at (-0.95, -1.6) {$n - \mathit{sub}(v)$ vertices,};
  \node[font=\sffamily\scriptsize, text=fgverm] at (-0.95, -1.95) {each 1 farther};
  \node[notel, anchor=west, text width=5.6cm] at (3.9, -1.0) {Moving the root from $p$ to $v$: every vertex in $v$'s subtree
    gets one step closer, all the others one step farther, so\\[3pt]
    \hspace*{1em}$\mathit{ans}(v) = \mathit{ans}(p) - \mathit{sub}(v) + (n - \mathit{sub}(v))$.};
\end{tikzpicture}
\caption{Rerooting the sum of distances across one edge.}
\label{fig:reroot}
\end{figure}

\begin{problem}{Tree Distances II (CSES)}
For every vertex, find the sum of its distances to all other vertices.
\end{problem}

\begin{code}
void solve() {
    int n;
    cin >> n;
    vvi adj(n + 1);
    rep(i, n - 1) {
        int a, b;
        cin >> a >> b;
        adj[a].push_back(b);
        adj[b].push_back(a);
    }
    // BFS order from vertex 1: parents always come before children
    vi order{1}, par(n + 1, 0), depth(n + 1, 0);
    rep(i, n) for (int v : adj[order[i]]) {
        if (v == par[order[i]]) continue;
        par[v] = order[i], depth[v] = depth[order[i]] + 1;
        order.push_back(v);
    }
    vi sub(n + 1, 1);  // subtree sizes, children first
    per(i, 1, n) sub[par[order[i]]] += sub[order[i]];
    vi ans(n + 1, 0);
    rep(v, 1, n + 1) ans[1] += depth[v];  // the sum of distances from 1
    // moving the root from p to its child v: sub[v] vertices get 1 closer,
    // the other n - sub[v] get 1 farther
    rep(i, 1, n) {
        int v = order[i];
        ans[v] = ans[par[v]] - sub[v] + (n - sub[v]);
    }
    print(vi(ans.begin() + 1, ans.end()));
}
\end{code}

\usetemplate{Rerooting (DP for every root)}{Ctrl+Alt+T, '}{cpreroot}

\subsection{When the combine cannot be undone}

The update above subtracts $v$'s contribution from $p$'s answer. That works for sums but not for maxima: knowing that the longest
path from $p$ has length 5 says nothing about the longest one once the branch into $v$ is removed. The standard fix is to keep the
\emph{two} best values over the children, so that the best value avoiding any single child is always at hand
(Figure~\ref{fig:reroot2}).

For the farthest distance from every vertex (Tree Distances~I again, now without the diameter argument), let
$\mathit{down}_1(v) \ge \mathit{down}_2(v)$ be the two longest downward paths from $v$ that start with different children, and let
$\mathit{up}(v)$ be the longest path from $v$ whose first step goes to its parent. For a child $v$ of $p$,
\[
  \mathit{up}(v) = 1 + \max\bigl(\mathit{up}(p),\ \text{the best downward path from $p$ that does not start with $v$}\bigr),
\]
where the second term is $\mathit{down}_2(p)$ if $\mathit{down}_1(p)$ goes through $v$, and $\mathit{down}_1(p)$ otherwise. The
farthest distance from $v$ is $\max(\mathit{down}_1(v), \mathit{up}(v))$: one pass children first for the downward values, one pass
parents first for the upward values.

\begin{figure}[htbp]
\centering
\begin{tikzpicture}
  \node[vtx] (g) at (0, 1.25) {};
  \node[vtx, hot] (p) at (0, 0) {$p$};
  \node[vtx, cool] (v) at (-1.7, -1.2) {$v$};
  \node[vtx] (c2) at (0, -1.2) {};
  \node[vtx] (c3) at (1.7, -1.2) {};
  \draw[ed, dashed] (g) -- ++(0, 0.75);
  \draw[ed] (g) -- (p) (p) -- (v) (p) -- (c2) (p) -- (c3);
  \foreach \nd/\hh in {v/1.5, c2/1.0, c3/0.7}
    \draw[draw=fgline, fill=fgfill, line width=0.6pt] (\nd.south) -- ++(-0.45, -\hh) -- ++(0.9, 0) -- cycle;
  \node[wl, text=fgblue] at ($(p)!0.5!(v)$) {5};
  \node[wl] at ($(p)!0.5!(c2)$) {3};
  \node[wl] at ($(p)!0.5!(c3)$) {2};
  \node[font=\sffamily\scriptsize, text=fgink, anchor=west] at ($(p)!0.5!(g)+(0.2, 0)$) {$\mathit{up}(p) = 3$};
  \node[font=\sffamily\scriptsize, text=fgblue] at (-1.7, -3.3) {$\mathit{down}_1(p) = 5$};
  \node[font=\sffamily\scriptsize, text=fgink] at (0, -2.8) {$\mathit{down}_2(p) = 3$};
  \node[notel, anchor=west, text width=6.2cm] at (2.9, -0.6) {The numbers on the edges are the longest downward paths from $p$
    through each child. The best one, 5, goes through $v$, so for $v$ the best path from $p$ that avoids $v$ is
    $\max(\mathit{up}(p), \mathit{down}_2(p)) = 3$ and $\mathit{up}(v) = 4$. For every other child it is
    $\max(\mathit{up}(p), \mathit{down}_1(p)) = 5$, giving $\mathit{up} = 6$.};
\end{tikzpicture}
\caption{Rerooting with a maximum: keeping the two best child values answers ``the best value without this child'' in $\Oh{1}$.}
\label{fig:reroot2}
\end{figure}

\begin{code}
void solve() {
    int n;
    cin >> n;
    vvi adj(n + 1);
    rep(i, n - 1) {
        int a, b;
        cin >> a >> b;
        adj[a].push_back(b);
        adj[b].push_back(a);
    }
    vi order{1}, par(n + 1, 0);
    rep(i, n) for (int v : adj[order[i]])
        if (v != par[order[i]]) par[v] = order[i], order.push_back(v);
    // down1[v] >= down2[v]: the two longest downward paths from v that start
    // with different children; up[v]: the longest path from v whose first
    // step goes to v's parent
    vi down1(n + 1, 0), down2(n + 1, 0), up(n + 1, 0);
    per(i, 1, n) {  // children first
        int v = order[i], p = par[v], len = down1[v] + 1;
        if (len > down1[p]) down2[p] = down1[p], down1[p] = len;
        else chmax(down2[p], len);
    }
    rep(i, 1, n) {  // parents first
        int v = order[i], p = par[v];
        // the longest path from p that does not go down into v
        int other = down1[p] == down1[v] + 1 ? down2[p] : down1[p];
        up[v] = 1 + max(up[p], other);
    }
    vi ans(n);
    rep(v, 1, n + 1) ans[v - 1] = max(down1[v], up[v]);
    print(ans);
}
\end{code}

When two children tie for the best value, $\mathit{down}_2(p) = \mathit{down}_1(p)$, so the test \cc{down1[p] == down1[v] + 1}
picks the right value either way.

\begin{trick}[Sum over all pairs: count what each edge separates]
The sum of the distances over all pairs of vertices needs no rerooting at all. The edge from $v$ to its parent lies on the path
between two vertices exactly when one of them is in $v$'s subtree and the other is not, so it contributes
$\mathit{sub}(v)\,(n - \mathit{sub}(v))$, and the total is $\sum_{v \ne r} \mathit{sub}(v)\,(n - \mathit{sub}(v))$, computed in
$\Oh{n}$. The same ``count the pairs an edge separates'' argument gives expected distances between random vertices, and, with the
edges sorted by weight and a DSU, the sum over all pairs of the heaviest edge on their path: an edge that joins components of
sizes $a$ and $b$ is the heaviest edge for exactly $a \cdot b$ pairs (Section~\ref{sec:contribution}).
\end{trick}

\begin{trick}[Designing a rerooting DP]
Write the answer for a fixed root as ``merge, over the children, of $\mathit{lift}$(child's value)''. If $\mathit{merge}$ is
associative with an identity, the generic template computes it for every root: the value of the ``rest of the tree'' seen from a
child is the merge of the parent's upward value and the siblings' values, which prefix and suffix merges provide without any
inverse. For the farthest vertex, $\mathit{merge} = \max$ and $\mathit{lift}(x) = x + 1$; for sums of distances, carry (sum, count)
pairs. Only invertible merges, like sums, allow the ``subtract the child'' shortcut used above.
\end{trick}

% ======================================================================================================================
\section[Merging child tables]{Merging child tables: the tree knapsack \hfill \plat}\label{sec:treeknap}

When the DP state contains a \emph{count} (``choose exactly $j$ vertices in the subtree''), merging a child into its parent is a
knapsack-style convolution, $\mathit{new}[a + b] = \max(\mathit{old}[a] + \mathit{child}[b])$. Done over full-size arrays this looks
like $\Oh{n^3}$. But if each table is only as long as its subtree, the total work is $\Oh{n^2}$.

\emph{Why.} Merging two tables of sizes $s_1$ and $s_2$ costs $s_1 s_2$, one unit for every pair of vertices with one vertex in each
part. Every pair of vertices is merged in this way exactly once, at their LCA, when their two branches are combined
(Figure~\ref{fig:treeknap}). There are $\binom{n}{2}$ pairs, hence $\Oh{n^2}$. If the tables are capped at $k$ entries, the same
counting gives $\Oh{nk}$.

\begin{figure}[htbp]
\centering
\begin{tikzpicture}
  \fill[fgsky!18, rounded corners=8pt] (-2.6, -1.55) rectangle (-0.35, -3.2);
  \fill[fgorange!15, rounded corners=8pt] (0.35, -1.55) rectangle (2.6, -3.2);
  \node[vtx, hot] (l) at (0, -0.6) {$v$};
  \node[vtx] (a) at (-1.5, -1.95) {1}; \node[vtx] (b) at (1.5, -1.95) {4};
  \node[vtx] (a1) at (-2.1, -2.8) {2}; \node[vtx, hot] (a2) at (-0.9, -2.8) {3};
  \node[vtx, hot] (b1) at (0.9, -2.8) {5}; \node[vtx] (b2) at (2.1, -2.8) {6};
  \draw[ed] (l) -- (a) (l) -- (b) (a) -- (a1) (a) -- (a2) (b) -- (b1) (b) -- (b2);
  \draw[line width=0.8pt, draw=fgverm, dashed] (a2) to[bend right=30] (b1);
  \node[font=\sffamily\scriptsize, text=fgverm] at (0, -3.6) {the pair $(3, 5)$ is paid for here, at $v$};
  \foreach \ii/\rw in {0/1, 1/2, 2/3} {
    \node[font=\sffamily\scriptsize\bfseries, text=fgblue] at (3.55, {-1.2 - 0.6*\ii}) {\rw};
    \foreach \jj in {0, 1, 2}
      \node[cellc, minimum width=6mm, minimum height=6mm] at ({4.0 + 0.6*\jj}, {-1.2 - 0.6*\ii}) {};
  }
  \foreach \jj/\cn in {0/4, 1/5, 2/6}
    \node[font=\sffamily\scriptsize\bfseries, text=fgorange!85!black] at ({4.0 + 0.6*\jj}, -0.68) {\cn};
  \node[cellc, minimum width=6mm, minimum height=6mm, fill=fgverm!25, draw=fgverm] at (4.6, -2.4) {};
  \node[lbl] at (4.6, -3.0) {merge loop: $3 \times 3$};
  \node[notel, anchor=west, text width=4.6cm] at (6.1, -1.8) {Merging the blue table into the orange one is a double loop over
    their sizes: one cell per pair of vertices, one from each side. The pair $(3, 5)$ is counted in this merge only, at its LCA
    $v$; over the whole run every pair is counted once.};
\end{tikzpicture}
\caption{Why merging subtree tables costs $\Oh{n^2}$ in total: the work of all merges together is at most one unit per pair of
vertices.}
\label{fig:treeknap}
\end{figure}

\begin{problem}{Choosing $k$ connected vertices}
Every vertex of a tree rooted at 1 has a weight (possibly negative). Choose exactly $k$ vertices such that the parent of every
chosen vertex other than the root is also chosen, with maximum total weight.
\solution $f_v[j]$ is the best weight of $j$ chosen vertices in $v$'s subtree forming such a closed set (so $v$ is chosen when
$j \ge 1$). Merge the children one by one with the knapsack convolution; a child may contribute only if $v$ itself is chosen.
\end{problem}

\begin{code}
void solve() {
    int n, k;
    cin >> n >> k;
    vi w(n + 1), par(n + 1, 0);
    rep(v, 1, n + 1) cin >> w[v];
    vvi child(n + 1);
    rep(v, 2, n + 1) {
        cin >> par[v];
        child[par[v]].push_back(v);
    }
    // f[v][j]: best weight choosing j vertices of v's subtree, closed under
    // parents (so v itself is chosen whenever j >= 1); arrays grow with size
    const ll NEG = LLONG_MIN / 4;
    vector<vi> f(n + 1);
    auto dfs = [&](auto&& self, int v) -> void {
        f[v] = {0, w[v]};
        for (int c : child[v]) {
            self(self, c);
            vi merged(min<int>(k, sz(f[v]) + sz(f[c]) - 2) + 1, NEG);
            rep(a, sz(f[v])) rep(b, sz(f[c])) {  // every pair of sizes once
                if (a + b >= sz(merged) || (a == 0 && b > 0)) continue;
                if (f[v][a] == NEG || f[c][b] == NEG) continue;
                chmax(merged[a + b], f[v][a] + f[c][b]);
            }
            f[v] = merged;
            f[c].clear();
        }
    };
    dfs(dfs, 1);
    print(k < sz(f[1]) ? f[1][k] : NEG);
}
\end{code}

\usetemplate{Tree knapsack (merging child DPs)}{Ctrl+Alt+Y, Shift+K}{cptreeknap}

Two details of the code matter. The merged table is only as long as the two tables together (capped at $k + 1$ entries), which is
what the counting argument needs; allocating $k + 1$ entries at every vertex would make every merge cost $\Theta(k^2)$, and
$\Theta(nk^2)$ in total. And a child's part may be non-empty only if $v$ itself is chosen ($a \ge 1$), which is how ``closed under
parents'' enters the merge.

\emph{Why the cap gives $\Oh{nk}$.} Call a part \emph{small} if it has fewer than $k$ vertices and \emph{big} otherwise. A merge of
a small part into a big one costs at most $k$ times the small part's size, and each vertex is on the small side of such a merge at
most once, since afterwards it belongs to a big part: $\Oh{nk}$ in total. Merges of two big parts cost $\Oh{k^2}$ each, and there
are fewer than $n / k$ of them, because each one reduces the number of disjoint big parts: $\Oh{nk}$ again. Merges of two small
parts are paid one unit per pair, and the pairs involved lie inside disjoint regions of fewer than $2k$ vertices each, so they
total at most $\sum 2k \cdot s_i = \Oh{nk}$.

% ======================================================================================================================
\section[Small-to-large merging]{Small-to-large merging \hfill \plat}\label{sec:smalltolarge}

To compute, for every vertex, something about the \emph{set} of values in its subtree (the number of distinct colours, the most
frequent value), give each vertex a set and merge the children's sets into the parent, always inserting the elements of the
\emph{smaller} set into the \emph{larger} one. An element is moved only when the set it belongs to at least doubles in size, so each
element moves $\Oh{\log n}$ times: $\Oh{n \log n}$ insertions in total, and $\Oh{n \log^2 n}$ time with \cc{std::set}. Swapping two
\cc{set}s is $\Oh{1}$, which is what makes ``keep the larger set'' free (Figure~\ref{fig:smalltolarge}).

\begin{figure}[htbp]
\centering
\begin{tikzpicture}
  \node[vtx, hot] (p) at (1.45, 0) {$p$};
  \node[draw=fgblue, fill=fgsky!18, rounded corners=6pt, minimum width=2.6cm, minimum height=0.8cm, line width=0.8pt]
    (L) at (0.3, -1.5) {};
  \node[draw=fgorange!85!black, fill=fgorange!15, rounded corners=6pt, minimum width=1.3cm, minimum height=0.8cm,
    line width=0.8pt] (S) at (3.0, -1.5) {};
  \foreach \xx in {-0.7, -0.2, 0.3, 0.8, 1.3} \fill[fgblue] (\xx, -1.5) circle (2.3pt);
  \foreach \xx in {2.75, 3.25} \fill[fgorange!85!black] (\xx, -1.5) circle (2.3pt);
  \draw[ed] (p) -- (L.north) (p) -- (S.north);
  \node[font=\sffamily\scriptsize\bfseries, text=fgblue, anchor=east] at (L.west) {larger: kept};
  \node[font=\sffamily\scriptsize\bfseries, text=fgorange!85!black, anchor=west] at (S.east) {smaller: moved};
  \draw[->, line width=1.1pt, draw=fgverm] (S.south) to[bend left=30]
    node[below, font=\sffamily\scriptsize, text=fgverm] {insert its 2 elements} (L.south);
  \foreach \kk/\ww/\tx in {0/0.3/start, 1/0.6/1 move, 2/1.2/2 moves, 3/2.4/3 moves} {
    \node[lbl, anchor=east] at (7.3, {0.1 - 0.55*\kk}) {\tx};
    \draw[draw=fgblue, fill=fgsky!18, line width=0.6pt, rounded corners=1.5pt] (7.45, {-0.1 - 0.55*\kk})
      rectangle ++(\ww, 0.4);
    \fill[fgverm] (7.6, {0.1 - 0.55*\kk}) circle (2.3pt);
    \node[lbl, anchor=west] at ({7.5 + \ww}, {0.1 - 0.55*\kk}) {size $\ge 2^{\kk}$};
  }
  \node[notel, anchor=north west, text width=4.9cm] at (6.0, -1.95) {The red element only ever moves into a set at least as
    large as its own, so each move at least doubles its set: at most $\log_2 n$ moves.};
\end{tikzpicture}
\caption{Small-to-large merging. Left: the parent keeps the larger child set (an $\Oh{1}$ swap) and inserts the elements of the
smaller one. Right: why each element moves $\Oh{\log n}$ times.}
\label{fig:smalltolarge}
\end{figure}

\begin{problem}{Distinct Colors (CSES)}
For every vertex of a tree, count the distinct colours in its subtree.
\end{problem}

\begin{code}
void solve() {
    int n;
    cin >> n;
    vi color(n + 1);
    rep(v, 1, n + 1) cin >> color[v];
    vvi adj(n + 1);
    rep(i, n - 1) {
        int a, b;
        cin >> a >> b;
        adj[a].push_back(b);
        adj[b].push_back(a);
    }
    vi order{1}, par(n + 1, 0);
    rep(i, n) for (int v : adj[order[i]])
        if (v != par[order[i]]) par[v] = order[i], order.push_back(v);
    vector<set<ll>> s(n + 1);
    vi ans(n + 1);
    per(i, n) {  // children first; then give the parent the bigger set
        int v = order[i];
        s[v].insert(color[v]);
        ans[v] = sz(s[v]);
        int p = par[v];
        if (!p) continue;
        if (sz(s[p]) < sz(s[v])) swap(s[p], s[v]);  // O(1) swap of sets
        for (ll x : s[v]) s[p].insert(x);  // insert the smaller into the larger
        s[v].clear();
    }
    print(vi(ans.begin() + 1, ans.end()));
}
\end{code}

\usetemplate{Small-to-large merging}{Ctrl+Alt+Y, M}{cpsmall2large}

The same pattern handles richer questions. For ``the most frequent colour in every subtree'', merge \cc{map<colour, count>}
tables small to large and keep, next to each table, its largest count and the sum of the colours that reach it; an insertion
updates both in $\Oh{1}$ after the map lookup. For queries ``how many vertices of $v$'s subtree lie at depth $d$'', merge
\cc{map<depth, count>}. The one rule is that the parent must \emph{adopt} the larger container, by \cc{swap} or by keeping an
index to it, and never copy it: copying the larger side costs its full size and destroys the bound.

\begin{trick}[DSU on tree (Sack)]
The same $\Oh{n \log n}$ bound holds with plain arrays instead of sets: process the heavy child (the one with the largest subtree)
last and keep its counts; then add the vertices of the light subtrees one by one, answer, and, if the current vertex is itself a
light child, erase everything on the way back. Each vertex is re-added once for every light edge above it, which is $\Oh{\log n}$
times (Chapter~\ref{ch:treedecomp}). It is faster and lighter on memory than merging sets, at the cost of more code.
\end{trick}

% ======================================================================================================================
\section[Tree isomorphism]{Tree isomorphism and hashing \hfill \plat}\label{sec:treeiso}

Two rooted trees are \emph{isomorphic} if one can be relabelled into the other while keeping the root, that is, if they have the
same shape up to reordering children. The \defn{AHU algorithm} decides it with canonical labels: process the vertices children first,
and give every vertex the label of the \emph{sorted} list of its children's labels, using a dictionary that assigns a new small
integer to every new list (Figure~\ref{fig:ahu}). Leaves get the label of the empty list. Two subtrees get the same label exactly when
they are isomorphic, by induction on height, so the trees are isomorphic exactly when the roots' labels are equal. With a shared
dictionary, all trees of a test file can be compared by their root labels, in $\Oh{n \log n}$ total.

\begin{figure}[htbp]
\centering
\begin{minipage}[c]{3.4cm}\centering
\begin{forest}
  cptree, for tree={s sep=6mm}
  [{3}, hot [{2}, cool [{0}] [{0}]] [{1}, warn [{0}]]]
\end{forest}
\end{minipage}
\hspace{6mm}
\begin{minipage}[c]{3.4cm}\centering
\begin{forest}
  cptree, for tree={s sep=6mm}
  [{3}, hot [{1}, warn [{0}]] [{2}, cool [{0}] [{0}]]]
\end{forest}
\end{minipage}
\hspace{6mm}
\begin{minipage}[c]{5cm}
\raggedright\small Labels: $() \to 0$, $(0) \to 1$, $(0, 0) \to 2$, $(1, 2) \to 3$. Both roots get the label of the sorted
list $(1, 2)$, so the trees are isomorphic although they are drawn differently.
\end{minipage}
\caption{AHU canonical labels. Each vertex shows the label of the sorted multiset of its children's labels.}
\label{fig:ahu}
\end{figure}

\begin{problem}{Tree Isomorphism I (CSES)}
Decide whether two trees rooted at vertex 1 are isomorphic.
\end{problem}

\begin{code}
void solve() {
    int t;
    cin >> t;
    map<vi, int> id;  // canonical id of each multiset of child ids (AHU)
    auto read_ids = [&](int n) {
        vvi adj(n + 1);
        rep(i, n - 1) {
            int a, b;
            cin >> a >> b;
            adj[a].push_back(b);
            adj[b].push_back(a);
        }
        vi order{1}, par(n + 1, 0), code(n + 1);
        rep(i, n) for (int v : adj[order[i]])
            if (v != par[order[i]]) par[v] = order[i], order.push_back(v);
        vvi kids(n + 1);
        per(i, n) {  // children first
            int v = order[i];
            sort(all(kids[v]));
            auto [it, fresh] = id.insert({kids[v], sz(id)});
            code[v] = it->second;
            if (par[v]) kids[par[v]].push_back(code[v]);
        }
        return code[1];  // equal ids <=> isomorphic rooted trees
    };
    while (t--) {
        int n;
        cin >> n;
        int a = read_ids(n), b = read_ids(n);
        yn(a == b);
    }
}
\end{code}

\emph{Complexity.} A vertex with $d$ children sorts $d$ labels and looks up a list of length $d$ in the dictionary, which costs
$\Oh{d \log n}$ with \cc{map<vi, int>}. The $d$ values sum to $n - 1$, so one tree costs $\Oh{n \log n}$. The labels are small
integers (at most the number of distinct subtree shapes seen so far), so they can index arrays directly afterwards.

\paragraph{Unrooted trees.} An isomorphism of unrooted trees maps centers to centers, and a tree has one center or two adjacent ones.
So root the first tree at one of its centers and compare its label with the second tree rooted at each of its (at most two) centers.
Centers are found by peeling leaves layer by layer: every round removes all current leaves, and the one or two vertices that remain
last are the centers, because peeling shortens every longest path by exactly 2.

\begin{problem}{Tree Isomorphism II (CSES)}
Decide whether two unrooted trees are isomorphic.
\end{problem}

\begin{code}
void solve() {
    int t;
    cin >> t;
    map<vi, int> id;
    auto rooted_id = [&](const vvi& adj, int root) {
        int n = sz(adj) - 1;
        vi order{root}, par(n + 1, 0), code(n + 1);
        rep(i, n) for (int v : adj[order[i]])
            if (v != par[order[i]]) par[v] = order[i], order.push_back(v);
        vvi kids(n + 1);
        per(i, n) {
            int v = order[i];
            sort(all(kids[v]));
            auto [it, fresh] = id.insert({kids[v], sz(id)});
            code[v] = it->second;
            if (par[v]) kids[par[v]].push_back(code[v]);
        }
        return code[root];
    };
    auto centers = [](const vvi& adj) {  // peel leaves layer by layer
        int n = sz(adj) - 1, left = n;
        vi deg(n + 1), layer;
        rep(v, 1, n + 1) {
            deg[v] = sz(adj[v]);
            if (deg[v] <= 1) layer.push_back(v);
        }
        while (left > 2) {
            left -= sz(layer);
            vi next;
            for (int v : layer)
                for (int u : adj[v])
                    if (--deg[u] == 1) next.push_back(u);
            layer = next;
        }
        return layer;  // one or two centers
    };
    while (t--) {
        int n;
        cin >> n;
        vvi g[2];
        rep(k, 2) {
            g[k].assign(n + 1, {});
            rep(i, n - 1) {
                int a, b;
                cin >> a >> b;
                g[k][a].push_back(b);
                g[k][b].push_back(a);
            }
        }
        // root the first tree at a center, compare with every center of the
        // second: an isomorphism maps centers to centers
        int a = rooted_id(g[0], centers(g[0])[0]);
        bool same = false;
        for (int c : centers(g[1])) same = same || rooted_id(g[1], c) == a;
        yn(same);
    }
}
\end{code}

\usetemplate{Tree isomorphism (canonical IDs, AHU)}{Ctrl+Alt+U, H}{cptreehash}

\begin{pitfall}[Hashing trees]
A popular shortcut hashes a vertex as, for example, $1 + \sum_c h(c)^2$ modulo a prime or with random polynomial evaluation. Such
hashes are fast and work well for counting distinct subtree shapes, but fixed, well-known hash functions have known collisions that
anti-hash tests exploit. Use the exact AHU labels when a wrong answer is fatal, or a hash with random parameters chosen at run time
(Section~\ref{sec:random-stl}).
\end{pitfall}

\begin{summary}[Chapter summary]
\begin{itemize}
  \item BFS order read backwards processes children before parents without recursion; preorder and postorder make subtrees
        contiguous; the full Euler tour turns LCA into a range minimum.
  \item Two BFS passes find a diameter, and the farthest vertex from anywhere is one of its ends.
  \item Euler tours turn subtree and root-path queries into range queries on a Fenwick tree.
  \item Binary lifting jumps $k$ levels and finds LCAs in $\Oh{\log n}$; jump pointers do it with one pointer per vertex; Tarjan's
        DSU algorithm answers LCAs offline in $\Oh{(n + q)\,\alpha(n)}$.
  \item Tree DP combines children bottom-up; rerooting gives every root's answer in two passes, subtracting a child for sums and
        keeping the two best children for maxima.
  \item Sums over all pairs of vertices are often sums over edges: each edge separates $\mathit{sub}(v)\,(n - \mathit{sub}(v))$
        pairs.
  \item Merging tables sized by subtrees costs $\Oh{n^2}$ in total (pairs meet once at their LCA); merging sets small to large costs
        $\Oh{n \log n}$ moves.
  \item AHU canonical labels decide rooted isomorphism; for unrooted trees, root at the centers.
\end{itemize}
\end{summary}

\begin{practice}
\ppgroup{Traversals, distances and queries}
\begin{pplist}
  \item Subordinates
  \item Tree Diameter
  \item Tree Distances I
  \item Tree Distances II
  \item Subtree Queries
  \item Path Queries
\end{pplist}
\ppgroup{Ancestors, DP and structure}
\begin{pplist}
  \item Company Queries I
  \item Company Queries II
  \item Distance Queries
  \item Tree Matching
  \item Distinct Colors
  \item Tree Isomorphism I
  \item Tree Isomorphism II
\end{pplist}
\end{practice}
